\section{Every original pulled-line global section retains its true norm on the same model}
\label{sec:native-pulled-projective-section-norms}

\begin{theorem}[The whole original section equivalence retains zero and nonzero sections {\cite[\S1.2]{BJ22}}]
\label{thm:native-pulled-projective-section-equivalence}
\lean{EnvelopeProperty.NativePulledProjectiveSectionNorms.coefficient_zero,EnvelopeProperty.NativePulledProjectiveSectionNorms.coefficient_ne_zero_iff,EnvelopeProperty.NativePulledProjectiveSectionNorms.sectionEquiv,EnvelopeProperty.NativePulledProjectiveSectionNorms.sectionEquiv_apply}
\leanok
\uses{thm:native-pulled-projective-model-sections,sc0-line-section-ratios:every-original-section}
Assume throughout this statement that the original scheme $X$ is integral.
For an actual original closed projective embedding, a nonempty original
pivot and any nonnegative degree $d$, the actual generic coefficient map
is an equivalence between the entire original global section set of
$f^*O(d)$ and the whole genuine degree-$d$ model section set. It sends zero
to zero and detects exactly the nonzero original global sections.
\end{theorem}
\begin{proof}
The preceding genuine coefficient membership, injection and actual
sheaf-gluing surjection give the equivalence, with its literal coefficient
as forward map. The canonical native section ratio is semilinear and
therefore sends zero to zero. Injectivity then identifies precisely the
nonzero original global sections with the nonzero model coefficients.
\end{proof}

\begin{theorem}[Every original global section has its genuine model graded section norm {\cite[(2.2)]{FHX}}]
\label{thm:native-pulled-projective-section-norms}
\lean{EnvelopeProperty.NativePulledProjectiveSectionNorms.sectionNorm,EnvelopeProperty.NativePulledProjectiveSectionNorms.sectionNorm_eq_chartValue,EnvelopeProperty.NativePulledProjectiveSectionNorms.sectionNorm_eq_secAbs}
\leanok
\uses{thm:native-pulled-projective-model-sections,thm:projective-native-model-global-section-values,thm:projective-native-model-value-one-center,thm:native-pulled-projective-line-frames,thm:fc-original-local-section-model-norm,lem:fc-original-local-section-scalar,lem:gl-sections}
Assume throughout this statement that the original scheme $X$ is integral.
At every original native analytic point, the true proper-model norm of
every actual original global section of the whole pulled $O(d)$ line equals
the genuine model graded section norm of its actual native coefficient,
under the same entire native analytic homeomorphism. No original ambient
homogeneous representation, generic-support assumption, chosen chart or
metric-comparison certificate is supplied.
\end{theorem}
\begin{proof}
The genuine model normalization supplies an actual coordinate of value
one. That coordinate is nonzero in the native function field, so its
original chart is nonempty. The proved whole support comparison puts the
original native support in that chart, while the genuine bounded-affine
criterion puts the original proper center there. In its actual original
pulled-line frame, the true original global-section model norm is evaluation
of the actual regular coefficient. The whole original chart equivalence
retains this coefficient's true value under the same global homeomorphism.
The preceding whole-section theorem identifies the corresponding chart
algebra element with the model fraction $g/D_j(d)$. The genuine model
section chart formula computes its graded norm as that fraction's chart
value times the original coordinate value to the degree $d$. The coordinate
value is one, so these two true norms coincide. The reasoning applies to
all original sections, including zero sections and native points with
nongeneric support; proper center and native support remain distinct.
\end{proof}
